%! AP Statistics — Unit 9, Lesson 9.2 — Student Packet Cover Sheet
\documentclass[10pt]{article}
\usepackage{apstats-article}
\usepackage{apstats-boxes}
\usepackage{ltablex}
\keepXColumns

\begin{document}

% Full-bleed navy banner
\begin{tikzpicture}[remember picture, overlay]
  \fill[navy] (current page.north west)
    rectangle ([yshift=-0.9in]current page.north east);
\end{tikzpicture}
\vspace{-0.8in}

\begin{center}
  {\color{white}\textbf{\LARGE AP Statistics}}\\[4pt]
  {\color{skymid}\large\textbf{Unit 9: Inference for Quantitative Data: Slopes}}\\[6pt]
  {\color{white}\textbf{\large Lesson 9.2 \quad Confidence Intervals for the Slope of a Regression Model}}
\end{center}

\vspace{0.22in}
\namedateperiod
\vspace{0.18in}

\begin{learningtargetbox}
\small
\begin{enumerate}[leftmargin=1.5em, itemsep=5pt]
  \item I can identify a $t$-interval for slope as the appropriate procedure for
        estimating the true slope $\beta$ of a regression model.
  \item I can verify the LINER conditions (Linear, Independent, Normal, Equal variance,
        Random) required for a confidence interval for slope.
  \item I can calculate a confidence interval for $\beta$ using $b \pm t^* \cdot SE_b$
        with $df = n-2$ and interpret it in context.
  \item I can read a computer regression output table and identify the slope $b$ and
        its standard error $SE_b$.
\end{enumerate}
\end{learningtargetbox}

\vspace{0.14in}

\begin{tocbox}
\small
\renewcommand{\arraystretch}{1.5}
\begin{tabularx}{\linewidth}{@{} c l X r @{}}
\rowcolor{sky}
\textbf{\#} & \textbf{Component} & \textbf{Description} & \textbf{Score} \\
\midrule
1 & Warm-Up        & Spiral review: LSRL, slope interpretation, residuals, meaning of $\beta=0$ & \blank{1.2cm} \\
2 & Guided Notes   & Sampling distribution of $b$; LINER conditions; $t$-interval formula; reading output & \blank{1.2cm} \\
3 & Group Activity & Fuel efficiency vs.\ weight: verify conditions, compute and interpret CI & \blank{1.2cm} \\
4 & Exit Ticket    & Independent: identify $b$ and $SE_b$, compute 95\% CI, interpret & \blank{1.2cm} \\
5 & Homework       & Two problems: full CI procedure and condition verification & \blank{1.2cm} \\
\midrule
  & \multicolumn{2}{r}{\textbf{Total}} & \blank{1.2cm} \\
\end{tabularx}
\end{tocbox}

\end{document}
